\documentclass[11pt,a4paper,fontset=none]{ctexart}
\usepackage[margin=18mm,top=18mm,bottom=18mm]{geometry}
\usepackage{fontspec,graphicx,booktabs,xcolor}
\setCJKmainfont{Noto Sans CJK SC}
\setmainfont{DejaVu Sans}
\definecolor{ink}{HTML}{183346}\definecolor{teal}{HTML}{087465}
\setlength{\parindent}{0pt}\setlength{\parskip}{5.5pt}\raggedbottom
\newcommand{\kicker}[1]{{\small\color{teal}\bfseries #1}\par}
\newcommand{\titleline}[1]{{\fontsize{20}{27}\selectfont\bfseries\color{ink}#1}\par\vspace{2mm}}
\newcommand{\topic}[1]{{\bfseries\color{teal}#1}\quad}
\newcommand{\chart}[2]{\begin{center}\includegraphics[width=\linewidth,height=#2\textheight,keepaspectratio]{#1.png}\end{center}}
\newcommand{\src}[1]{\vfill{\footnotesize\color{gray}证据入口\quad #1}\par}
\begin{document}\color{ink}

\kicker{实验决策报告　／　2026 年 11 月 20 日}

\titleline{局部码字能解释什么，重建误差又来自哪里？}

阶段 A 冻结三个已有硬码模型，检验单个两秒码字能解释的局部属性。阶段 B 在相同训练预算下改变编码器输入与量化条件。两组结果分别回答模式含义与重建瓶颈，评价对象均为开发验证被试。

\begin{center}\small\begin{tabular}{lrrrr}
\toprule
开发验证宏 NMSE & X 连续 & X 硬码 & 目标连续 & 目标硬码 \\
\midrule
EEG & 0.22957 & 0.38694 & 0.20180 & 0.37308 \\
Hb & 0.15498 & 0.18514 & 0.00703 & 0.04574 \\
\bottomrule\end{tabular}\end{center}

\topic{EEG}连续模型直接读取目标后，宏误差从 0.22957 降到 0.20180，但 REFED 总电压 30–45 Hz 功率仍比目标低 41.86 dB。两种输入下硬量化又分别增加 0.15737 和 0.17127 的宏误差。频带保真与量化适配均有缺口。

\topic{Hb}目标输入连续模型的宏误差为 0.00703，远低于 X 输入的 0.15498。G4 的局部 ID 对双色团斜率更有区分力，同时重建和曲率读出仍有取舍。

\topic{下一步}建议立即开展一项 EEG 连续模型的频带平衡对照。首轮保留 E-TC 的目标输入、结构、训练抽样和 10,000 步预算，只改变失真权重。Hb 的训练扩展和 K512 码本待这次判别后再决定。

\topic{阅读顺序}第 2–8 页给出设计、成对差值、频谱、波形和收敛轨迹。第 9–13 页给出单码属性、真实同码样本、跨库误差及换窗结果。第 14 页列明实验选择与判定条件。

\src{阶段 A 三个模型终态、阶段 B 八个匹配单元的固定终点。附表 AB\_tables.xlsx 提供各库、被试、码字的数值。}

\newpage

\kicker{设计与核验}

\titleline{两组实验共享开发人口，各有独立验收指标}

\topic{目标与接口}观测信号 X 分解为既有生理模型（teacher）估计的条件分量 S 与剩余量 O = X − S。EEG 使用 200 Hz 电压，Hb 使用 10 Hz 的 HbO、HbR。每个分支在完整 30 秒上下文中每两秒输出一个位置，共 15 个。K128 每分支每位置一个整数，名义码率 3.5 bit/s；解码器只接收本分支整数序列。完整序列利用了双向上下文。G1 为普通波形解码器；G4 的 Hb 解码器用当前码决定局部形状，限制序列上下文的形状修正。

\topic{输入与预热}实际应用路径从 X 预测 S、O，两分支相加得到 X 的重建。特权路径给各分支其自身目标，用来诊断有限预算的压缩误差。EEG 均使用八个片内位置的 E1 结构，Hb 使用 G1 waveform。E-Q 与 H-TQ 复用 E1、D1 的预热，E-TC 与 E-TQ 从同一个新预热分叉。每组 1,000 步预热加 9,000 步延续，seed 17，位置选择和训练尺度固定。

\begin{center}\small\begin{tabular}{lrrrrr}
\toprule
验证人口 & Single-Trial & REFED & Simultaneous & Visual & 合计 \\
\midrule
被试 & 6 & 6 & 5 & 3 & 20 \\
\bottomrule\end{tabular}\end{center}

\topic{分母}每模态 144 条记录。EEG 有 1,920 个窗口、65,909 个窗口位置输入，Hb 有 1,923 个窗口、70,972 个输入。G0 已保留六个 EEG 上游失败父窗口。阶段 A 按完整两秒片段计分，EEG 为 955,640 / 988,635，Hb 为 1,034,497 / 1,064,580；阶段 B 在各窗口有效时间点上计分。

\topic{误差}各 X、S、O 分量和各色团分别以有效样本平方误差均值除以同库同位置冻结训练目标方差。由位置到窗口、记录、被试、数据集逐层等权，再对分量和色团等权得宏 NMSE。值越低越好。水平误差项采用每个两秒片段的平均残差平方，形状误差项是去掉该均值后的残差平方，两项相加等于 NMSE。

\topic{核验}阶段 A 三组审计均完成；阶段 B 四个新增模型均完成，各有 4k、7k、10k 检查点。与四个历史控制组逐模态核对训练评分、验证和训练人口内回归面板的最终输入身份，并从逐输入 NMSE 重新汇总验证被试值与宏值。固定 10k 为主终点。

\src{G0 人口和尺度合同、A/B resolved\_config、plan、各模型 summary 及 evaluation/step\_10000。}

\newpage

\kicker{阶段 B　／　配对差值}

\titleline{两种模态对输入与量化的敏感程度不同}

\chart{contrasts}{.28}

\begin{center}\small\begin{tabular}{lrrrr}
\toprule
模态 & 比较 & 差值 & 95\% 配对区间 & 正差 / 负差被试 \\
\midrule
EEG & 连续输入差 & +0.02777 & [+0.02299, +0.03328] & 20/0 \\
EEG & X 量化差 & +0.15737 & [+0.13168, +0.18295] & 20/0 \\
EEG & 目标量化差 & +0.17127 & [+0.14363, +0.19890] & 20/0 \\
EEG & 交互差 & -0.01390 & [-0.01674, -0.01125] & 0/20 \\
Hb & 连续输入差 & +0.14795 & [+0.13319, +0.16524] & 20/0 \\
Hb & X 量化差 & +0.03016 & [+0.02006, +0.04784] & 20/0 \\
Hb & 目标量化差 & +0.03871 & [+0.02529, +0.06261] & 20/0 \\
Hb & 交互差 & -0.00855 & [-0.01491, -0.00467] & 1/19 \\
\bottomrule\end{tabular}\end{center}

连续输入差取 X 连续减目标连续；另两项在相同输入下取硬码减连续。交互差取 X 量化差减目标量化差。正值表示相减的第一项误差更大。EEG 的 X 量化差约为连续输入差的 5.7 倍，Hb 的连续输入差约为 X 量化差的 4.9 倍。

固定训练种子、模型与四个数据库，库内按被试有放回配对重采样 10,000 次，每次四库等权，取 2.5\% 与 97.5\% 分位数。区间描述这组冻结模型在开发人口上的不确定性，未作多重比较校正。两模态交互均为负，数据没有显示量化对 X 输入施加额外的更大惩罚。

\src{八个 10k 验证输出的逐被试 NMSE；附表 B\_差值。}

\newpage

\kicker{阶段 B　／　EEG}

\titleline{EEG 的输入改善主要落在 S，硬码影响四库}

\chart{EEG_matrix}{.49}

每格均为验证 NMSE，颜色越深误差越大。上图每格汇总 X、S、O，下图每格对四库等权；两幅图色标分别标出。输入改善最大的是连续 S，从 0.3989 降至 0.3180。连续总电压 X 仅由 0.1416 变为 0.1414，说明宏误差改善不能直接读作总电压波形恢复。

在 X 输入下，连续模型的宏水平 / 形状 NMSE 为 0.0082 / 0.2214，硬码为 0.0219 / 0.3651。目标输入下两项分别为 0.0001 / 0.2017 和 0.0133 / 0.3598。硬码代价的主要贡献来自片内形状。

换回历史四位置 G1 硬码作补充比较，E-Q 的宏 NMSE 从 0.39752 降至 0.38694，只有 8/20 名被试下降。四库中 REFED 改善，另三库的均值略升。REFED 的配对均值差为 −0.04853，另三库为 +0.00072 至 +0.00334。这个结果支持保留 E-Q 为匹配基线，优先解决频带缺口。

\src{阶段 B 固定验证终点；历史 G1 waveform K128 的同身份验证面板。}

\newpage

\kicker{阶段 B　／　EEG 频谱}

\titleline{直接看到目标后，REFED 高频仍明显衰减}

\chart{spectrum}{.5}

各曲线对应同一库总电压 X。对 30 秒窗口在有效点去均值，无效点置零，乘 Hann 窗后计算频带功率比，先取对数再按评分层级汇总，乘 10 得 dB。各带左闭右开。零线表示预测与目标功率相同，负值表示衰减。图中各库共用纵轴。

REFED 目标连续组在 1–4 Hz 为 −0.21 dB，在 13–30 和 30–45 Hz 为 −33.86 和 −41.86 dB。X 连续在后两带为 −38.80、−45.22 dB。直接给定分支目标改善了部分高频恢复，剩余衰减仍很大。四库目标连续模型在 30–45 Hz 的偏差均落在 −42.46 至 −37.38 dB。

目标硬码在后两带达到 −53.17、−52.24 dB。量化增加了频带损失，但连续路径本身已显著压低高频功率。频带结果支持先在连续路径检验损失权重；功率与相位保真要分别检查。

\src{evaluation/step\_10000/validation 的 X 分量频带对数功率比。}

\newpage

\kicker{阶段 B　／　EEG 固定样例}

\titleline{同一两秒电压片段上，硬码的波形误差更大}

\chart{EEG_example}{.48}

\begin{center}\small\begin{tabular}{lrrr}
\toprule
示例 RMSE & X & S & O \\
\midrule
X 连续 & 0.405 & 0.217 & 0.286 \\
X 硬码 & 0.696 & 0.293 & 0.483 \\
目标连续 & 0.411 & 0.214 & 0.281 \\
目标硬码 & 0.721 & 0.295 & 0.489 \\
\bottomrule\end{tabular}\end{center}

REFED 被试 17，video\_1，M1，窗口 0–30 秒，展示 14–16 秒。对保存的 REFED 固定示例按原顺序取中间位置，选择独立于误差。黑线分别为观测 X、teacher 的 S、O；彩线为四组同窗预测。整窗可见，纵轴均以训练 X 标准差为单位，各行按分量取纵轴。

X 连续与目标连续在该片段的总电压误差接近，分别为 0.405 与 0.411。两种硬码总电压误差均在 0.69 以上，O 分量从连续的约 0.28 升到硬码的约 0.48。示例的 RMSE 使用上图所示时间点，验证频谱结论来自上一页的完整面板。

\src{四组 validation\_examples 的共同 row/start/position，检查目标和掩码一致。}

\newpage

\kicker{阶段 B　／　Hb}

\titleline{Hb 的主要差距集中在从观测估计剩余分量}

\chart{Hb_matrix}{.51}

各行 X、S、O 中的 HbO/HbR 分别为氧合与脱氧血红蛋白。两张图分别在库内汇总和四库等权汇总，色标各自给出。连续 X 输入的 HbO-O / HbR-O 误差为 0.4333 / 0.3056，直接给定目标后降至 0.0038 / 0.0042；X 自身的误差则一直较小。原路径能够贴近总曲线，同时允许两个分量出现较大的相反残差。

X 输入连续组宏水平 / 形状 NMSE 为 0.1407 / 0.0143；目标连续为 0.0059 / 0.0012。H-TQ 为 0.0344 / 0.0114，整体仍优于 X 输入连续。四个库均保留该排序。特权输入测出当前预算下的分解学习差距，无法据此把全部差距归为 X 缺少信息。

在总曲线 X 上，H-TQ 的 HbO/HbR 相关为 0.951/0.952，X 输入硬码为 0.960/0.970。H-TQ 的宏分量优势与总曲线的相关性改善是两个不同的结果。

\src{四个 Hb 单元的最终 NMSE、水平/形状分解及相关。}

\newpage

\kicker{阶段 B　／　Hb 固定样例与训练轨迹}

\titleline{Hb 分量误差缩小，四个新增模型仍在改善}

\chart{Hb_example}{.34}

\begin{center}\small\begin{tabular}{lrrr}
\toprule
示例 RMSE & X & S & O \\
\midrule
X 连续 & 0.050 & 0.125 & 0.141 \\
X 硬码 & 0.069 & 0.119 & 0.147 \\
目标连续 & 0.008 & 0.024 & 0.020 \\
目标硬码 & 0.043 & 0.037 & 0.045 \\
\bottomrule\end{tabular}\end{center}

REFED 被试 17，video\_1，CH26，展示同一完整 30 秒的 HbO。样例按 EEG 页的相同规则选择。S 分量的目标硬码 RMSE 为 0.037，X 输入硬码为 0.119；两条分量在总量内相加，单看 X 会低估二者的差异。HbR 的全体结果见上一页与附表。

\chart{trajectory}{.23}

折线为 4k、7k、10k 的完整验证宏 NMSE。新增四组的最小已观测值均在 10k。H-TQ 从 7k 的 0.05129 降到 0.04574。这些差值反映等预算表现，训练收敛尚无更长轨迹验证。

\src{固定 HbO 示例及所有八个配置的三个保留检查点。}

\newpage

\kicker{阶段 A　／　EEG 单码与序列}

\titleline{当前 ID 可读出部分电压形态，高频属性增益弱}

\chart{EEG_readout}{.42}

同一完整两秒片段上，背景读出使用库、原生位置与局部观测 RMS 档（边界 0.5、1、2 个训练 X SD）。前三种读出使用训练片段条件均值；完整序列用同分支冻结解码器恢复波形，再提取相同属性。上图为 100 × (1 − 背景加 ID 误差 / 背景误差)；下图以单 ID 误差作分母。正值表示新增条件后误差降低。

未经对齐的 EEG 原电压在 G1-S/O 加入 ID 后分别改善 21.6\% / 45.7\%，去均值形状为 17.4\% / 17.8\%。两分支 30–45 Hz 的增益分别为 −1.8\% / −0.8\%。O 分支 1–4 Hz 宏增益为 −99.5\%；Simultaneous 为 +33.6\%，另三库为负。该属性的条件均值向验证人口迁移存在明显库间差异。

五个频带为两秒去均值 Hann 谱的积分功率，包络为八个 250 ms 箱的 RMS。频带属性 MSE 评价带内功率，波形 MSE 保留实际相位。下图比较两种已经拟合的读出，有限解码器的负增益只说明该读出在该属性上误差更大。这里评价的 G1 与阶段 B 的 E-Q 拥有不同前端。

\src{A/G1\_eeg/summary.json。工作表 A\_读出含背景、单 ID、联合及序列的绝对 MSE。}

\newpage

\kicker{阶段 A　／　Hb 单码与序列}

\titleline{G4 将更多斜率信息集中到当前码字}

\chart{Hb_readout}{.44}

列内依次表示同色团的原波形、两秒均值、线性斜率及二次拟合曲率。斜率为每秒变化率，曲率为每秒平方的二阶导；每个片段 20 个时间点。增益分母和符号与上一页相同，四种读出对同一批完整目标评分。

HbO 斜率在 S 分支的背景上 ID 增益从 G1 的 26.0\% 升到 G4 的 61.9\%，在 O 分支从 4.9\% 升到 47.9\%。HbR 斜率也从 23.2\% / 44.9\% 升到 56.2\% / 58.0\%。按同一被试的背景加 ID 斜率 MSE 配对，HbO-S 与 HbO-O 均有 20/20 名被试下降，HbR-S/O 分别为 17/20、18/20。与此同时，S 分支 HbO 曲率增益从 15.3\% 降到 9.6\%。

以 HbO-S 斜率为例，背景、单 ID、背景加 ID、序列 MSE 在 G1 为 0.0680、0.0505、0.0503、0.0197，在 G4 为 0.0680、0.0247、0.0259、0.0229。G4 的优势集中在局部读出；序列波形重建仍有代价。历史 H1 匹配解码器的 G1/G4 宏 NMSE 为 0.18265/0.20880。

\src{A/G1\_hb 与 G4\_hb，历史 D0/D1/H1 的冻结赋码匹配解码器结果。}

\newpage

\kicker{阶段 A　／　逐码支持与跨库迁移}

\titleline{128 个码均有训练支持，跨库覆盖差异明显}

\chart{support}{.32}

\chart{transfer}{.32}

上图每点一个码，下图每个盒子汇总一个库内有验证样本的码。盒内横线为中位数，盒边为四分位数，须长至 1.5 倍四分位距内的最远点，圆点为其外的码。下图在库内以训练码字均值预测验证波形，RMSE 以训练 X SD 表示；Hb 展示 HbO。纵轴取对数，各面板独立标度。ST/REF/SIM/VIS 依次对应四个数据库。

六个分支各 128 码均达到训练至少 3 名被试、20 个完整片段。四库均有验证支持的码数依次为 G1 EEG S/O 97/74、G1 Hb 122/127、G4 Hb 128/128。G1 EEG-O 最少的码只有 21 个验证片段、3 名被试。四个库均参与训练，这里的迁移评价针对各库留出的被试。

\src{各分支 atlas，附表提供全部 768 个分支码及分库条件方差、验证原型误差。}

\newpage

\kicker{阶段 A　／　真实同码样本}

\titleline{同一码字仍覆盖具有不同幅度和形状的片段}

\chart{gallery}{.44}

每个分支选择完整验证片段最多的码。各取 D0 已保存的 64 个同码真实片段，逐段减去自身均值；细灰线为其中 3 个真实片段，深色线为逐时刻中位数，阴影为 10\%–90\% 分位带。图中 EEG 为原相位电压，Hb 为 HbO，各图独立纵轴；该带描述样本分布。

例如，G1 EEG-O 的 ID 91 覆盖 91,795 个完整验证片段，G4 Hb-S 的 ID 55 覆盖 48,149 个。波形带只画保留样本。用全部验证片段逐码计算时，先逐时间点求条件方差，再对时间平均并开方。128 个码等权的中位离散度在 G1 EEG-S/O 为 0.185/0.706 个训练 X SD；G1 HbO-S/O 为 0.643/0.337，G4 为 0.569/0.345。

上述离散度使用未去均值的波形，因此包含水平变化。配合上一页的支持、迁移误差及单码属性读出，G4 Hb-S 支持可重复的局部方向属性；给单个码指定唯一精细波形仍超出这些结果。

\src{D0 对应 validation\_same\_ID 的真实片段；阶段 A 全量 atlas 的条件方差。}

\newpage

\kicker{阶段 A　／　自然上下文}

\titleline{G4 换码更少，发生换码时原型差异反而更大}

\chart{context}{.32}

在同一验证父窗口、同一位置内，从两个相邻且准入的 30 秒评分块构造真正向后平移 2 秒的输入。比较原输入第 8 个与平移后第 7 个两秒片段，两者物理时间一致。EEG 有 34,629 对，Hb 有 37,052 对；分别有 31,280、33,920 个位置输入缺少所需后继块。每个配对的中央片段均完整。

左图按位置、窗口、记录、被试与库逐层汇总。G1 EEG-S/O 换码率为 5.23\% / 6.24\%；G1 Hb 为 27.05\% / 35.91\%，G4 Hb 为 10.52\% / 19.64\%。原窗冻结模型 CPU 复算与历史 ID 的宏不一致率最高 0.17\%。

右图只对切换且有训练支持的配对求两个训练原型波形的均方差，再对片段池化平均并开方。HbO-S 在 G1/G4 为 0.302/0.533，HbO-O 为 0.173/0.282 个训练 X SD。因此换码频率与换码时的属性变化须分别评价。

训练形态重编号后六个分支的 128 个有支持码均能零总代价匹配。当前只有 seed 17；独立训练种子之间的形态复现，以及独立生理事件关联，仍缺少直接证据。

\src{各分支 natural\_context 配对数组与训练原型；计数和条件均方变化已从配对记录复核。}

\newpage

\kicker{决策　／　一个可判别的下一步}

\titleline{立即做一个 EEG 连续频带平衡对照}

\topic{优先级依据}目标连续模型在 REFED 仍损失近 42 dB 的 30–45 Hz 功率，且连续总电压的片内形状 NMSE 与 X 输入几乎相同。调整这一连续模型的频带失真，可以直接检验训练目标是否足以解释部分高频衰减。直接扩大词表只能检查量化环节。

\topic{匹配设计}建议以 E-TC 为控制，在同一 G0 训练人口、E1 八位置前端、seed 17、抽样及 10k 预算下，对 S、O、X 引入预先固定的非均匀频带权重，保留复数频谱相位与现有时域、差分评价。权重及各频带尺度仅由训练集确定，预先设置低能量频带的尺度下限；均匀加权的复数频谱平方误差与时域平方误差等价。若沿用原预热，则从同一 1k checkpoint 分叉、只在后续更新改变损失。这是本报告建议的新探索设计。

\topic{继续条件}运行前固定主要比较及可接受的时域代价。按四库、同被试检验 REFED 13–30 与 30–45 Hz 的功率偏差是否同时向零靠近；并列核对原相位时域 X/S/O NMSE、波形相关与其他库变化。仅当频带改善且既定时域容差满足时，才把这一改变带回可部署的 X 输入连续组，随后检查整数路径。

\topic{Hb 与局部模式}Hb 的连续输入差为 0.14795，约为 X 路径量化差的 4.9 倍。保留 G1 作为重建对照，保留 G4 作为斜率模式参照。若独立安排 Hb 的下一项诊断，应比较 X 输入连续模型有界延长训练后的 O 分量误差，并保持 teacher 和评分面板冻结。本轮优先只开展上述一项 EEG 对照。

\topic{停止与分流}若提升频带功率伴随原相位波形明显变差，应保留该负结果；若连续频带已改善而硬量化仍带来较大损失，再单独检验 K512。针对可部署候选，可在确定候选及匹配基线后安排两个额外种子。当前结论适用于四库开发验证与单种子等预算诊断。

\topic{可复核材料}附表给出四个库及每名被试的重建、背景与 ID 读出、每码支持和条件方差、换窗结果、训练轨迹与区间。对应运行目录为 20261120\_pattern\_a\_v1 与 20261120\_input\_quantization\_b\_v1。此次复核未启动后续训练。

\src{阶段 A/B 运行合同、固定终点与阶段 A 审计；实验建议以本次开发结果为依据。}

\end{document}
